\chapter{Vector Instructions and Data-Parallel Code}\label{ch:ll:vector-instructions-and-data-parallel-code}

A parser that reads one byte at a time looks at a message the way a person reads a page letter by letter. In 2019 Geoff
Langdale and Daniel Lemire published a JSON parser that looks at 32 or 64 bytes per instruction instead, and described it as
the first standard-compliant JSON parser ``to process gigabytes of data per second on a single core'', using a quarter or
fewer of the instructions of a widely used reference parser. On the laptop of this book, the loop that looks for the end of a
field one byte at a time scans about four bytes a nanosecond; the same search written with 256-bit \omterm{def:ll:vector-instructions-and-data-parallel-code:simd}{vector instructions} scans
about eighty, some eighteen times faster, and splits a 159-byte FIX order into its sixteen fields in about \qty{11}{\nano\second}
instead of a hundred. Market data and order messages are text or packed binary, read field by field on every tick: this
chapter shows what \omterm{def:ll:vector-instructions-and-data-parallel-code:simd}{vector instructions} are, when the compiler uses them by itself and what stops it, how to write them by
hand with a scalar fallback that the tests hold them to, and how to lay out data so that they apply.

\section{SIMD registers and instruction sets}

\begin{definition}[SIMD, vector instruction]\label{def:ll:vector-instructions-and-data-parallel-code:simd}
\emph{SIMD}\index{SIMD} (single instruction, multiple data) is the execution of one operation on several data elements at
once. A \emph{vector instruction}\index{vector instruction} is such an operation on a vector register holding several
elements side by side: on x86-64, 16-byte \texttt{xmm} registers (SSE2, part of every x86-64 processor), 32-byte
\texttt{ymm} registers (AVX2) and, on processors that have it, 64-byte \texttt{zmm} registers (AVX-512).
\end{definition}

A 32-byte register holds 32 bytes, 16 two-byte integers, 8 four-byte integers or floats, or 4 doubles, and one instruction
adds, compares or multiplies all of them. The instruction that matters most for parsing is the byte comparison:
\texttt{VPCMPEQB} compares the 32 bytes of two \texttt{ymm} registers and sets each byte of the result to all ones where they
are equal, and \texttt{VPMOVMSKB} gathers the top bit of each of those bytes into a 32-bit integer in an ordinary register. A
message chunk compared with a register full of the delimiter byte therefore becomes a 32-bit mask with one bit per position,
and ``where is the first delimiter'' becomes ``how many trailing zeros does the mask have'', one more instruction
(\cref{fig:ll:vector-instructions-and-data-parallel-code:mask}). The loop does in three instructions what the scalar loop
does in 32 comparisons and 32 branches.

\begin{omfigure}
\begin{tikzpicture}[font=\scriptsize,c/.style={draw=gray!60,minimum width=0.34cm,minimum height=0.42cm,inner sep=0pt}]
\foreach \ch [count=\j from 0] in {3,5,=,D,|,5,5,=,E,S,Z,6,|,5,4,=,1,|,3,8,=,5,|,4,4,=,5,7,2,3,|,6} {
  \pgfmathtruncatemacro{\i}{\j}
  \pgfmathtruncatemacro{\hit}{(\i==4)||(\i==12)||(\i==17)||(\i==22)||(\i==30)}
  \ifnum\hit=1
    \node[c,fill=omThm!25] at (0.36*\i,0) {\ch};
    \node[c,fill=omThm!55,text=white] at (0.36*\i,-1.3) {1};
  \else
    \node[c,fill=gray!8] at (0.36*\i,0) {\ch};
    \node[c] at (0.36*\i,-1.3) {0};
  \fi
  \node[c,fill=omDef!12] at (0.36*\i,-0.65) {|};
}
\node[anchor=east,font=\footnotesize] at (-0.3,0) {32 bytes of a message};
\node[anchor=east,font=\footnotesize] at (-0.3,-0.65) {delimiter, broadcast};
\node[anchor=east,font=\footnotesize] at (-0.3,-1.3) {\texttt{movemask} bits};
\foreach \i in {0,8,16,24,31} {\node[font=\tiny,gray] at (0.36*\i,0.38) {\i};}
\draw[-Stealth,thick,omThm] (0.36*4,-1.55) -- ++(0,-0.45) node[below,align=center,font=\footnotesize]
  {trailing zeros $= 4$:\\ first delimiter};
\node[anchor=west,align=left,font=\footnotesize] at (3.4,-2.25) {clear the lowest bit and count again: 12, 17, 22, 30\\
  (one \texttt{VPCMPEQB}, one \texttt{VPMOVMSKB}, one count per field)};
\end{tikzpicture}

\omcaption{Finding the field delimiters of a FIX message with one comparison of 32 bytes. The byte \texttt{|} stands for the
SOH delimiter (byte 1). The comparison sets a byte to all ones where the message equals the delimiter; the mask
instruction collects one bit per byte; counting trailing zeros gives each position in
turn.}\label{fig:ll:vector-instructions-and-data-parallel-code:mask}
\end{omfigure}

The laptop's processor has SSE2 through AVX2 and fused multiply-add (One Quant Book~4, chapter~25), but not AVX-512, as its
feature flags in \path{/proc/cpuinfo} show. That is the ordinary situation for a trading firm's code: it is compiled once
and runs on machines of several generations, so the widest instruction set cannot be assumed, and a binary that executes an
AVX2 instruction on a processor without it dies with an illegal-instruction signal. Section~3 deals with that.

\section{Auto-vectorisation and what blocks it}

\begin{definition}[Auto-vectorisation]\label{def:ll:vector-instructions-and-data-parallel-code:auto}
\emph{Auto-vectorisation}\index{auto-vectorisation} is the compiler's transformation of a scalar loop into \omterm{def:ll:vector-instructions-and-data-parallel-code:simd}{vector
instructions} that process several iterations at once, done when it can prove the result unchanged and estimates it
profitable.
\end{definition}

The compiler vectorises by itself more often than people expect, and less often than they hope. GCC does it at
\texttt{-O3}; from GCC~12 it also does it at \texttt{-O2}, with a cost model that accepts only very cheap
transformations (the release notes), so a loop that GCC~11 at \texttt{-O2} leaves scalar may be vectorised by GCC~13 at the
same flag: the compiler version is part of the performance. \texttt{-fopt-info-vec-all} makes it report each decision.
\Cref{lst:ll:vector-instructions-and-data-parallel-code:report} is its report on five small loops under the laptop's GCC~11.

\omcode{code/low-latency/14-vector-instructions-and-data-parallel-code/cpp/vecreport/gxx11.txt}{1}{20}{The vectorisation
report of GCC~11 for \texttt{ll\_vecdemo.cpp}: the revaluation over columns (line 13), over records (20--21), a search that
leaves early (27), and integer and double sums (34, 40).}\label{lst:ll:vector-instructions-and-data-parallel-code:report}

Each line is a lesson. At \texttt{-O2} nothing is vectorised. The revaluation over separate arrays is vectorised at
\texttt{-O3}, two doubles per instruction, four with \texttt{-mavx2}, but only after the compiler ``versioned'' the loop:
the output pointer might overlap an input, so it checks at run time and keeps a scalar copy of the loop for that case
(\texttt{\_\_restrict}, which promises no overlap, removes the check). The same computation over 80-byte records is not
vectorised at all: the five fields of consecutive options are 80 bytes apart, and GCC~11 finds no vector type for such
strided loads. The search is refused because it may leave the loop early (``control flow in loop''): GCC~11 does not vectorise a loop
with an early exit. The integer sum is vectorised. The double sum is reported
vectorised too, and its assembly shows why that means nothing: the additions stay in their original order, one element at a
time, because floating-point addition is not associative and reordering it would change the result (chapter~8, where
\texttt{-ffast-math} allowed it and changed the answer). Blockers, in short: possible aliasing, strided or scattered data,
early exits, calls that cannot be inlined, and reductions whose order matters.

\section{Intrinsics and CPU feature dispatch}

\begin{definition}[Compiler intrinsic, CPU feature dispatch]\label{def:ll:vector-instructions-and-data-parallel-code:intrinsic}
A \emph{compiler intrinsic}\index{compiler intrinsic} is a function the compiler knows and translates into one specific
instruction (\texttt{\_mm256\_cmpeq\_epi8} into \texttt{VPCMPEQB}), letting a program use an instruction without writing
assembly while the compiler still allocates registers and schedules around it. \emph{CPU feature
dispatch}\index{CPU feature dispatch} selects at run time, from the features the processor reports, which of several
compiled versions of a function to call.
\end{definition}

When the compiler will not vectorise a loop, intrinsics do it by hand. The build's search for a byte loads 32 bytes, compares
them with the broadcast delimiter, takes the mask and counts its trailing zeros; a message whose length is not a multiple of
32 ends with one more load that finishes exactly at its last byte and overlaps the previous block, with the bytes already
seen shifted out of the mask, so that no byte beyond the message is ever read. That last point is not pedantry: a read past
the end of a buffer that ends at a page boundary faults, and the sanitiser of chapter~8 reports any such read.

\omcode{code/firm/simdscan/cpp/firm_simdscan.hpp}{64}{79}{The AVX2 search for a byte, with an overlapping last
block.}\label{lst:ll:vector-instructions-and-data-parallel-code:find}

The function is compiled for AVX2 by a \texttt{target} attribute, without compiling the whole program for it, and is called
only after a check. GCC's \texttt{\_\_builtin\_cpu\_supports(\textquotedbl{}avx2\textquotedbl{})} returns non-zero when the running processor has the
feature; the build checks once, at the first call, and keeps the answer. Chapter~8's \omterm{def:ll:cpp-for-latency-iii-the-compiler:fmv}{function multiversioning} does the same
automatically for a whole function; the explicit version keeps the choice visible and testable, and the tests call every
version directly on the same inputs.

\omcode{code/firm/simdscan/cpp/firm_simdscan.hpp}{196}{210}{Dispatch: detect the features once, then call the widest
version.}\label{lst:ll:vector-instructions-and-data-parallel-code:dispatch}

Rust expresses the same contract in its types. An intrinsic-using function is declared \texttt{unsafe} with the attribute \texttt{target\_feature} enabling \texttt{avx2}, and calling it is \omterm{def:ll:rust-for-low-latency:sound}{sound}
only once the macro \texttt{is\_x86\_feature\_detected!} has returned true for \texttt{avx2}, as the standard library's own example puts it: the
\texttt{unsafe} block ``is safe because we're testing that the \texttt{avx2} feature is indeed available on our CPU''.

\omcode{code/firm/simdscan/rust/src/lib.rs}{86}{98}{The same dispatch in Rust: the unsafe call is guarded by the run-time
check.}\label{lst:ll:vector-instructions-and-data-parallel-code:rust}

Numbers need a second trick. Eight ASCII digits loaded as one 64-bit integer can be turned into their value without a loop:
subtract the character \texttt{0} from every byte at once, then combine neighbouring digits into two-digit numbers in 16-bit
lanes, those into four-digit numbers in 32-bit lanes, and those into the result, three multiplications in all. This is \omterm{def:ll:vector-instructions-and-data-parallel-code:simd}{SIMD}
within an ordinary register, and it needs no special instruction; a version with SSE4.1 multiply-adds does the same in vector
registers.

\omcode{code/firm/simdscan/cpp/firm_simdscan.hpp}{134}{144}{Eight digits in three multiply-add steps inside one 64-bit
register.}\label{lst:ll:vector-instructions-and-data-parallel-code:swar}

\begin{proposition}[The eight-digit combination]\label{prop:ll:vector-instructions-and-data-parallel-code:swar}
Let bytes $b_0,\dots,b_7$ hold digits $d_0,\dots,d_7$ (after subtracting the character zero), $b_0$ the lowest. After the
step $v \leftarrow (10v + (v \gg 8)) \wedge \mathtt{0x00FF00FF00FF00FF}$, the 16-bit lane $k$ holds $10 d_{2k} + d_{2k+1}$;
after the two analogous steps with 100 and 10\,000, the low 32 bits hold $\sum_{j=0}^{7} d_j\, 10^{7-j}$.
\end{proposition}

\begin{proof}
Before the mask, byte $2k$ of $10v + (v \gg 8)$ is $10 d_{2k} + d_{2k+1} \le 99 < 256$, since the shift brings byte $2k+1$
down to position $2k$; no carry leaves the byte, and the mask discards the odd bytes, which hold unwanted sums. The same
argument on 16-bit lanes (values up to $99 \times 100 + 99 < 2^{16}$) and on 32-bit lanes (up to $9\,999 \times 10\,000 +
9\,999 < 2^{32}$) gives the four-digit and eight-digit values, the earlier digit always in the higher place.
\end{proof}

On the laptop an eight-digit field costs about \qty{2.0}{\nano\second} with the scalar loop, \qty{1.0}{\nano\second} with the
register trick and \qty{0.7}{\nano\second} with SSE4.1, against about \qty{5}{\nano\second} for the standard library's
\texttt{std::from\_chars}, which also validates and handles any length. A price parsed into an integer number of
ten-thousandths costs about \qty{7}{\nano\second}, and the same price read as a double by \texttt{strtod} about
\qty{43}{\nano\second}. Lemire's number-parsing paper makes the point for floating-point parsing: a 64-bit float needs at most
17 significant digits, which fit in one 64-bit word, and parsing done that way is several times faster than the C library's.

\section{Where vectorisation pays in a trading system}

\Cref{fig:ll:vector-instructions-and-data-parallel-code:scan} measures the search for a delimiter placed at the end of a
message of each length. The scalar loop costs about a quarter of a nanosecond per byte; SSE2 about a fortieth, AVX2 about an
eightieth, on top of a fraction of a nanosecond per call. Below 16 bytes both vector versions hand the message to the scalar loop
and cost the same. Splitting a whole FIX order message into its fields (16 fields, 159 bytes) takes about
\qty{95}{\nano\second} one byte at a time and about \qty{11}{\nano\second} with the mask loop, about nine times less.

\begin{omfigure}
\begin{tikzpicture}
\begin{axis}[width=11.5cm,height=5.8cm,xmode=log,ymode=log,log basis x=2,xmin=6,xmax=5500,ymin=0.4,ymax=2000,
  xlabel={message length (bytes, log scale)},ylabel={ns per search (log scale)},
  xtick={8,16,32,64,128,256,512,1024,4096},xticklabels={8,16,32,64,128,256,512,1024,4096},
  tick label style={font=\small},label style={font=\small},grid=major,grid style={gray!25},
  legend image code/.code={\draw[#1] (0cm,0cm) -- (0.6cm,0cm);},
  legend columns=3,legend style={font=\footnotesize,at={(0.5,-0.3)},anchor=north,draw=none,fill=none,
  /tikz/every even column/.append style={column sep=8pt}},table/col sep=comma]
\addplot[omThm,very thick,mark=*,mark size=1.2pt] table[x=bytes,y=scalar]{figdata/low-latency/14-vector-instructions-and-data-parallel-code/measured_scan.csv};
\addplot[omMeth,thick,mark=square*,mark size=1.2pt] table[x=bytes,y=sse2]{figdata/low-latency/14-vector-instructions-and-data-parallel-code/measured_scan.csv};
\addplot[omDef,very thick,mark=triangle*,mark size=1.6pt] table[x=bytes,y=avx2]{figdata/low-latency/14-vector-instructions-and-data-parallel-code/measured_scan.csv};
\legend{scalar loop,SSE2 (16 bytes),AVX2 (32 bytes)}
\end{axis}
\end{tikzpicture}

\omcaption{Time to find a delimiter at the end of a message, by message length: median of batches of a thousand searches of
a message in the first-level cache. Measured on a laptop (Intel Core Ultra 7 155H) under WSL2, no isolated cores. Data:
\texttt{bench\_simd.py}.}\label{fig:ll:vector-instructions-and-data-parallel-code:scan}
\end{omfigure}

Where does this matter? Wherever bytes are examined one by one on the hot path: splitting FIX messages on their delimiter
(chapter~15), finding the structural characters of JSON messages from venues that publish in it (chapter~17), validating and
parsing numeric fields, checksums, and copying. It matters less for binary protocols whose fields sit at fixed offsets
(chapter~16), where there is nothing to search for; there, \omterm{def:ll:vector-instructions-and-data-parallel-code:simd}{vector instructions} pay in bulk work on arrays instead: the risk
checks of many orders, the revaluation of a book of positions, the statistics a strategy updates. Three conditions decide:
enough elements to fill a register, the same operation on each, and data laid out contiguously. The first two are properties
of the problem; the third is a design decision.

\section{Structure of arrays}

\begin{definition}[Structure of arrays]\label{def:ll:vector-instructions-and-data-parallel-code:soa}
A \emph{structure of arrays}\index{structure of arrays} stores a collection of records as one array per field, so that the
same field of consecutive records is contiguous in memory; the usual layout, one record after another, is called an
array of structures.
\end{definition}

An options book revalued from its Greeks shows the difference
(\cref{fig:ll:vector-instructions-and-data-parallel-code:layout}). Each option's record carries an identifier, a symbol, its
quantity and four Greeks, an expiry and some flags: 80 bytes, of which the revaluation $q(\Delta\,\delta S + \tfrac12
\Gamma\,\delta S^2 + \mathcal{V}\,\delta\sigma + \Theta\,\delta t)$ reads 40. Stored as records, four consecutive deltas lie
80 bytes apart and no instruction loads them together; stored as columns, they are 32 consecutive bytes, one load. Half of
every \omterm{def:ll:memory-hierarchy-and-caches:line}{cache line} fetched for the records is wasted as well.

\begin{omfigure}
\begin{tikzpicture}[font=\scriptsize]
\node[anchor=west,font=\footnotesize\bfseries] at (-0.1,1.0) {array of structures (80-byte records)};
\foreach \r in {0,1,2} {
  \pgfmathsetmacro{\x}{\r*4.3}
  \draw[fill=gray!10] (\x,0) rectangle ++(0.8,0.5) node[midway]{id};
  \draw[fill=gray!10] (\x+0.8,0) rectangle ++(1.1,0.5) node[midway]{symbol};
  \foreach \f/\k in {q/0,$\Delta$/1,$\Gamma$/2,$\mathcal{V}$/3,$\Theta$/4} {
    \draw[fill=omDef!25] (\x+1.9+0.4*\k,0) rectangle ++(0.4,0.5) node[midway]{\f};
  }
  \draw[fill=gray!10] (\x+3.9,0) rectangle ++(0.35,0.5);
}
\foreach \r in {0,1,2} {\draw[-Stealth,omThm,thick] (\r*4.3+2.5,-0.05) -- (\r*4.3+2.5,-0.45);}
\node[anchor=west,omThm] at (-0.1,-0.7) {deltas 80 bytes apart: one load each, half of every line unused};
\node[anchor=west,font=\footnotesize\bfseries] at (-0.1,-1.45) {structure of arrays (one array per field)};
\foreach \f/\y in {$\Delta$/-2.15,$\Gamma$/-2.75} {
  \node[anchor=east] at (-0.15,\y+0.25) {\f};
  \foreach \k in {0,...,15} {\draw[fill=omDef!25] (0.4*\k,\y) rectangle ++(0.4,0.5);}
}
\draw[very thick,omThm] (0,-2.15) rectangle ++(1.6,0.5);
\node[anchor=west,omThm] at (6.6,-1.9) {four deltas, 32 contiguous bytes:};
\node[anchor=west,omThm] at (6.6,-2.3) {one \texttt{ymm} load, one fused multiply-add};
\end{tikzpicture}

\omcaption{The same options book as records and as columns. Revaluation reads five of the record's fields; as columns, each
field of four consecutive options fills one 256-bit register.}\label{fig:ll:vector-instructions-and-data-parallel-code:layout}
\end{omfigure}

\Cref{fig:ll:vector-instructions-and-data-parallel-code:reval} measures both layouts at three flag sets. With 1\,024
options, which fit in the caches, the records cost about \qty{0.8}{\nano\second} per option whatever the flags, since no
compiler vectorises them; the columns cost the same at \texttt{-O2} (GCC~11 does not vectorise there), half at \texttt{-O3}
(two doubles per instruction) and a quarter with \texttt{-mavx2 -mfma}, four doubles per instruction and one fused
multiply-add for each multiplication and addition pair. With a million options the data no longer fit in the caches and the
bytes moved decide: 88 bytes per option for records (the record and the result), 48 for columns, and the measured cost is
about three times lower for columns whatever the flags, the vector width now mattering little.

\begin{omfigure}
\begin{tikzpicture}
\begin{axis}[width=12cm,height=5.8cm,ybar,bar width=8pt,ymin=0,ymax=6.5,
  symbolic x coords={O2,O3,O3avx2},xtick=data,xticklabels={\texttt{-O2},\texttt{-O3},\texttt{-O3 -mavx2 -mfma}},
  enlarge x limits=0.25,ylabel={ns per option},tick label style={font=\small},label style={font=\small},grid=major,
  grid style={gray!25},legend columns=2,legend style={font=\footnotesize,at={(0.5,-0.22)},anchor=north,draw=none,fill=none,
  /tikz/every even column/.append style={column sep=8pt}},table/col sep=comma]
\addplot[fill=gray!35,draw=gray] table[x=key,y=rows_1k]{figdata/low-latency/14-vector-instructions-and-data-parallel-code/measured_reval.csv};
\addplot[fill=omDef!70,draw=omDef] table[x=key,y=cols_1k]{figdata/low-latency/14-vector-instructions-and-data-parallel-code/measured_reval.csv};
\addplot[fill=gray!70,draw=gray!80!black] table[x=key,y=rows_1m]{figdata/low-latency/14-vector-instructions-and-data-parallel-code/measured_reval.csv};
\addplot[fill=omDef!35,draw=omDef] table[x=key,y=cols_1m]{figdata/low-latency/14-vector-instructions-and-data-parallel-code/measured_reval.csv};
\legend{{records, 1\,024 options},{columns, 1\,024 options},{records, $2^{20}$ options},{columns, $2^{20}$ options}}
\end{axis}
\end{tikzpicture}

\omcaption{Revaluing an options book from its Greeks, stored as 80-byte records or as columns, compiled at three flag sets by
GCC~11: in cache (1\,024 options) and from memory ($2^{20}$ options, 80~MB of records). Measured on a laptop (Intel Core
Ultra 7 155H) under WSL2, no isolated cores. Data: \texttt{bench\_simd.py}.}\label{fig:ll:vector-instructions-and-data-parallel-code:reval}
\end{omfigure}

Columns have costs. Adding or removing one option touches every array; code that handles one option at a time (an order
arriving, a fill) reads one element from each of several arrays, which is more \omterm{def:ll:memory-hierarchy-and-caches:line}{cache lines} than one record; and the
layout is less natural to write. The usual compromise is to keep records where single items are handled and columns where
many items are computed on, or blocks of columns (eight options' deltas, then their gammas) that give both. A book of orders
handled one message at a time (chapter~19) wants records; a risk engine revaluing every position on every move (chapter~22)
wants columns.

\section{Tutorial: scan, parse and lay out}\label{tut:ll:vector-instructions-and-data-parallel-code:tut}

\begin{tutorial}
\textbf{Goal.} Scan and parse with \omterm{def:ll:vector-instructions-and-data-parallel-code:simd}{vector instructions} behind a feature check, held byte for byte to a scalar reference; read
the compiler's decisions; measure layouts. \textbf{End state:}
\cref{fig:ll:vector-instructions-and-data-parallel-code:scan,fig:ll:vector-instructions-and-data-parallel-code:reval},
\cref{lst:ll:vector-instructions-and-data-parallel-code:report}, and green tests in C++ and Rust.
\begin{tutsteps}
\item \textbf{Hold the vector code to the scalar code.} The build's C++ test runs every length from 0 to 200 bytes with
delimiters at random positions (and none), twenty messages per length, through the scalar, SSE2, AVX2 and dispatched
searches and the position finders, and requires identical answers; it parses 200\,000 eight-digit numbers four ways. The
buffers are exactly as long as the messages, so a run under AddressSanitizer catches any read past the end.
\item \textbf{Read the compiler's report}: \texttt{python ll\_vecreport.py} compiles \texttt{ll\_vecdemo.cpp} at three flag
sets and writes \cref{lst:ll:vector-instructions-and-data-parallel-code:report} for the installed GCC; run it with a newer
GCC and compare the \texttt{-O2} lines.
\item \textbf{Measure} with \texttt{python bench\_simd.py}: searches by length, the split of FIX messages, the parsers, and
the revaluation built at three flag sets.
\item \textbf{The same in Rust}: \texttt{cargo test} in \path{code/firm/simdscan/rust} runs the twin's scans for every length
and its parsers on the same values.
\end{tutsteps}
\textbf{What to change next.} Add \texttt{\_\_restrict} to the pointers of \texttt{reval\_columns} and check that the
report's versioning line disappears; write the double sum with four separate accumulators and see whether the compiler
vectorises it for real.
\end{tutorial}

\section{Build: the firm's scanner}\label{bld:ll:vector-instructions-and-data-parallel-code:simdscan}

\begin{build}
\bfield{Purpose} The byte-level primitives of every text protocol the firm reads: the FIX engine splits messages on the
delimiter with it (chapter~15), and the WebSocket client finds JSON's structural characters and parses prices and sizes with
it (chapter~17).
\bfield{Interface} C++20 \texttt{firm::simdscan}: \texttt{find\_byte(p, n, c)}, \texttt{positions(p, n, c, out, cap)},
\texttt{positions\_any(p, n, set, m, out, cap)}, each with \texttt{\_scalar}, \texttt{\_sse2} (search) and \texttt{\_avx2}
versions; \texttt{parse8\_scalar}, \texttt{parse8\_swar}, \texttt{parse8\_sse}, \texttt{is8digits},
\texttt{parse\_uint(p, n, value, used)}, \texttt{parse\_fixed(p, n, scale, out)}; \texttt{level()}. Rust
\texttt{firm\_simdscan}: \texttt{find\_byte}, \texttt{positions}, \texttt{parse8\_swar}, \texttt{parse\_uint},
\texttt{parse\_fixed}, with the scalar references.
\bfield{Rules} Every vector routine returns what its scalar reference returns for every input, never reads outside
$[p, p+n)$, and is called only after the run-time feature check; the AVX2 code is compiled by function attributes, never by
compiling the program for AVX2; prices are parsed into integers of a declared scale, never into doubles.
\bfield{Acceptance tests} \texttt{code/firm/simdscan/}: every length from 0 to 200 bytes against the scalar references in
C++ and Rust, the capacity limit of the position finders respected, 200\,000 eight-digit values parsed identically by every
method, integers of 19 digits, prices with and without decimals, negative prices, and the rejection of malformed ones.
\bfield{Stretch} An AVX-512 version selected by the same dispatch; a classifier for JSON's structural characters with a
table lookup (one shuffle instead of one comparison per character); a validator for eight digits inside the vector path.
\end{build}

\begin{omsources}
\item G.~Langdale and D.~Lemire, ``Parsing gigabytes of JSON per second'', \emph{The VLDB Journal} 28(6), 2019.
\item D.~Lemire, ``Number parsing at a gigabyte per second'', \emph{Software: Practice and Experience} 51(8), 2021.
\item GCC 12 release notes (vectorisation at \texttt{-O2}); GCC manual, x86 built-in functions
(\texttt{\_\_builtin\_cpu\_supports}).
\item Intel 64 and IA-32 Architectures Software Developer's Manual, instructions \texttt{PCMPEQB} and \texttt{PMOVMSKB} (as
transcribed in the x86 instruction reference at felixcloutier.com).
\item The Rust standard library, \texttt{core::arch} and \texttt{is\_x86\_feature\_detected}.
\end{omsources}

\section{Exercises}

\begin{exercise}[$\star$]\label{exo:ll:vector-instructions-and-data-parallel-code:1}
How many doubles, 32-bit integers and bytes does a 256-bit register hold? How many 256-bit comparisons does it take to look
at every byte of a 1\,500-byte packet?
\end{exercise}

\begin{exercise}[$\star$]\label{exo:ll:vector-instructions-and-data-parallel-code:2}
The mask of a 32-byte block is \texttt{0x00041010} (bit 0 the first byte). Where are the delimiters in the block?
\end{exercise}

\begin{exercise}[$\star$]\label{exo:ll:vector-instructions-and-data-parallel-code:3}
Apply \cref{prop:ll:vector-instructions-and-data-parallel-code:swar} by hand to the digits \texttt{20260925}: give the
16-bit lanes after the first step and the 32-bit lanes after the second.
\end{exercise}

\begin{exercise}[$\star\star$]\label{exo:ll:vector-instructions-and-data-parallel-code:4}
A message of 70 bytes is scanned with 32-byte blocks and an overlapping last block. Which bytes does each load cover, how
many bits must be shifted out of the last mask, and why is the result still correct if the delimiter is at byte 60?
\end{exercise}

\begin{exercise}[$\star\star$]\label{exo:ll:vector-instructions-and-data-parallel-code:5}
Explain each of the three reasons the compiler gives in \cref{lst:ll:vector-instructions-and-data-parallel-code:report}
for not vectorising, or for vectorising only after a check, and how you would remove each.
\end{exercise}

\begin{exercise}[$\star\star$]\label{exo:ll:vector-instructions-and-data-parallel-code:6}
From the measured scan, estimate how long AVX2 takes to search the whole of a 4\,096-byte buffer and what that is in bytes
per nanosecond. What limits it once the buffer no longer fits in the first-level cache?
\end{exercise}

\begin{exercise}[$\star\star\star$]\label{exo:ll:vector-instructions-and-data-parallel-code:7}
\emph{Coding.} Add \texttt{positions\_any\_sse2} to \texttt{firm.simdscan} (16-byte blocks, overlapping last block), use it
in the dispatch when AVX2 is absent, and extend the C++ test to compare it with the scalar reference for every length.
\end{exercise}

\begin{exercise}[$\star\star\star$]\label{exo:ll:vector-instructions-and-data-parallel-code:8}
\emph{Find the flaw.} ``Our build machine has AVX-512, so we compile with \texttt{-march=native}; it passed every test in
continuous integration and ran fine on the new servers.''
\end{exercise}

\section{Problem: Thirty-Two Bytes at a Time}

\begin{problem}[{Weekend problem --- when does a vector scan pay?}]\label{pb:ll:vector-instructions-and-data-parallel-code:1}
The measured scan (\texttt{measured\_scan.csv}) gives the cost of each search by message length; the split and parser
measurements are in \texttt{measured\_split.csv} and \texttt{measured\_parse.csv}. A FIX order gateway receives execution
reports of about 200 bytes with 20 fields, 200\,000 a second at peak.

\medskip\noindent\textbf{Part I --- The costs.}
\begin{enumerate}
\item From the measurements at 4\,096 bytes, what does each method cost per byte?
\item Fit a line $a + bL$ to the AVX2 costs for $L \ge 32$ (\texttt{ll\_simd.fit}). What are $a$ and $b$?
\item Charging the scalar loop only its cost per byte, above what length is the AVX2 search cheaper? And SSE2?
\item Why is charging the scalar loop no fixed cost generous to it?
\end{enumerate}

\medskip\noindent\textbf{Part II --- The gateway.}
\begin{enumerate}[resume]
\item What does splitting one 200-byte report cost with each method, scaled from the measured split?
\item How much core time per second does splitting take at peak, each way?
\item Each report carries four prices and quantities to parse. Using \texttt{parse\_fixed} rather than \texttt{strtod}, what
is saved per report?
\item Which is the larger saving per report: the split or the parsing?
\end{enumerate}

\medskip\noindent\textbf{Part III --- The limits.}
\begin{enumerate}[resume]
\item Why does the vector search not help a binary protocol with fields at fixed offsets?
\item What happens on a processor without AVX2?
\item Why must the last block overlap rather than read the 32 bytes beyond the start of the tail?
\item What would AVX-512 change in the cost per byte, and what would it not change?
\end{enumerate}

\medskip\noindent\textbf{Part IV --- The verdict.}
\begin{enumerate}[resume]
\item State the \emph{named result}: the message length below which the vector scan does not repay its setup, from the
measured cost per byte and per call.
\item Is any FIX field shorter than that?
\item What measurement would you add before trusting the break-even on a server?
\item Why are the scalar references kept in the build?
\item How would you test that the dispatch picks the scalar path on a machine without AVX2, on a machine that has it?
\item Where else in a trading system would the same thirty-two-bytes-at-a-time pattern apply?
\item What layout change would let the risk checks of chapter~22 use the same instructions?
\item In one sentence: when do \omterm{def:ll:vector-instructions-and-data-parallel-code:simd}{vector instructions} pay?
\end{enumerate}
\end{problem}

\section{Interview questions}

\begin{interviewq}[$\star$ \iqroles{developer}]\label{iq:ll:vector-instructions-and-data-parallel-code:1}
What is \omterm{def:ll:vector-instructions-and-data-parallel-code:simd}{SIMD}? Give an example of a loop that benefits and one that does not.
\end{interviewq}

\begin{interviewq}[$\star\star$ \iqroles{developer}]\label{iq:ll:vector-instructions-and-data-parallel-code:2}
How would you find the first occurrence of a byte in a buffer as fast as possible?
\end{interviewq}

\begin{interviewq}[$\star\star$ \iqroles{developer}]\label{iq:ll:vector-instructions-and-data-parallel-code:3}
What stops a compiler from vectorising a loop? How do you find out why it did not?
\end{interviewq}

\begin{interviewq}[$\star\star$ \iqroles{developer}]\label{iq:ll:vector-instructions-and-data-parallel-code:4}
Your binary uses AVX2 and must run on older machines. How do you ship it?
\end{interviewq}

\begin{interviewq}[$\star\star$ \iqroles{developer, researcher}]\label{iq:ll:vector-instructions-and-data-parallel-code:5}
Array of structures or \omterm{def:ll:vector-instructions-and-data-parallel-code:soa}{structure of arrays} for a position book that is revalued on every price move? Why?
\end{interviewq}

\begin{interviewq}[$\star\star\star$ \iqroles{developer}]\label{iq:ll:vector-instructions-and-data-parallel-code:6}
Parse an ASCII decimal price into a fixed-point integer without a loop over characters. What must you validate?
\end{interviewq}
